\documentclass[12pt]{article}

\usepackage[utf8]{inputenc}
\usepackage[T1]{fontenc}
\usepackage{geometry}
\geometry{a4paper, margin=1in}
\usepackage{hyperref}
\usepackage{enumitem}
\usepackage{titling}

\title{DNAx Lab \\ \large Codebase Analysis and Technical Report}
\author{Autogenerated Technical Report}
\date{\today}

\begin{document}

\maketitle

\begin{abstract}
This report provides a comprehensive summary of the DNAx Lab codebase, detailing its architecture, main components, features, and underlying structure. The application is a desktop tool designed to assist with synthetic biology and molecular biology workflows, featuring a modern graphical user interface and robust in-silico validation tools.
\end{abstract}

\section{Architecture}
The DNAx Lab application is a desktop application written in Python using the \texttt{tkinter} library for the Graphical User Interface. It is structured using a clear component-based pattern where the UI pages (Views) handle user interaction and delegate complex logic to modular backend components.

\begin{itemize}
    \item \textbf{Entry Point (\texttt{src/main.py}):} Initializes the main application window (\texttt{App} class inheriting from \texttt{tk.Tk}), sets up the window properties (icon, maximized state), and builds the main UI, providing a collapsible side navigation menu and a main content area.
    \item \textbf{Navigation (\texttt{src/navigation.py}):} Provides a \texttt{NavigationFrame} container that holds multiple page frames and manages switching between them seamlessly.
    \item \textbf{UI Widgets (\texttt{src/ui\_widgets.py}):} Contains custom, modern-looking \texttt{tkinter} widgets, such as \texttt{ModernButton}, \texttt{NavButton}, \texttt{PlaceholderEntry}, and an animated \texttt{CircularLoader}.
    \item \textbf{Utilities (\texttt{src/utils.py}):} Includes helper functions like \texttt{resource\_path} for resolving asset paths (crucial for PyInstaller packaging) and uninstaller hooks.
\end{itemize}

\section{Main Features \& Tools}
The application is modularized into several specialized tools and pages dedicated to distinct scientific calculations and designs:

\subsection{Calculators and Generators}
\begin{itemize}
    \item \textbf{Size Calculator (\texttt{src/tools/size\_calc.py}):} Calculates the linear length, approximate circular diameter, estimated molecular weight (in Da and kDa), and average mass per base given the number of bases. It assumes a typical B-form DNA rise of 0.34 nm per base.
    \item \textbf{DNA Generator (\texttt{src/tools/dna\_generate.py}):} Generates random DNA sequences adhering to specific constraints (GC content 40--60\%, no long homopolymers, avoidance of specific restriction sites like BsaI). It supports both Circular (Golden Gate Assembly) and Linear Fragment modes, accompanied by advanced plasmid map plotting.
\end{itemize}

\subsection{Design and Analysis Tools}
\begin{itemize}
    \item \textbf{Primer Designer (\texttt{src/tools/primer\_designer.py}):} Designs optimal PCR primer pairs considering Tm, GC\%, 3' stability, and avoidance of hairpins and dimers. It also supports designing cloning primers that include restriction enzyme overhangs and clamp sequences.
    \item \textbf{qPCR Analysis \& Probe Design (\texttt{src/tools/qpcr.py}):} Finds optimal TaqMan/Hydrolysis probes based on Tm, GC content, and sequence rules. Includes an Efficiency Calculator and a Relative Quantification ($\Delta\Delta$Ct) calculator to determine fold changes.
    \item \textbf{BLAST Comparator (\texttt{src/tools/comparator.py}):} Interfaces with the NCBI BLAST API to search for sequence matches in the nucleotide database, categorizing hits into Genomic/Natural and Synthetic/Vector matches.
\end{itemize}

\subsection{Validation and Documentation}
\begin{itemize}
    \item \textbf{In-Silico Validation (\texttt{src/pages/validation.py}):} Acts as a comprehensive QC report. It fetches the current design from the DNA Generator, validates PCR compatibility, verifies digestion sites, and checks sticky-end compatibility for ligation. Visualizations are provided by a custom \texttt{PlasmidMap}.
    \item \textbf{Export (\texttt{src/pages/export.py}):} Allows exporting the generated DNA construct details (sequence, primers, probes, summary) to Excel (.xlsx), PDF, and Cloud (Firebase).
    \item \textbf{Protocols (\texttt{src/pages/protocol.py}):} A built-in PDF viewer using \texttt{PyMuPDF} to display common lab protocols stored locally in the \texttt{assets/protocols} directory.
\end{itemize}

\section{Supporting Modules and Tests}
\begin{itemize}
    \item \textbf{Simulation Logic (\texttt{src/tools/simulation.py}):} Contains \texttt{run\_simulation}, the backend for the In-Silico Validation page, parsing restriction sites, sticky ends, and checking internal cut validity.
    \item \textbf{Automated Tests:} The root directory contains \texttt{test\_dna\_generator.py} (a script simulating the 3-phase DNA generation algorithm) and \texttt{test\_simulation.py} (unit tests for the assembly logic).
\end{itemize}

\section{UI/UX and Dependencies}
The application features a modern, flat design aesthetic within Tkinter, utilizing custom hover effects and rounded shapes drawn on standard canvas elements. Essential dependencies include the Python standard library (\texttt{tkinter}, \texttt{math}, \texttt{threading}), as well as external packages such as \texttt{requests} (for API interactions), \texttt{openpyxl} (for Excel export), \texttt{reportlab} \& \texttt{qrcode} (for PDF creation), and \texttt{PyMuPDF} \& \texttt{Pillow} (for the embedded PDF viewer).

\end{document}
